\documentclass[publication-type=technical-report,theme=default]{reportkit}
% Quiz fixture: true/false (circle one) and multiple choice.
% Student sheet. quiz-key.tex builds the same source as an answer key.
% All content below is invented for this fixture.
\usepackage{reportkit-quiz}
\setreportkitleftheader{Unit 3 quiz}
\setreportkitfooter{Fictional course fixture}
\setreportkitversion{Fixture 1.0}
\setreportkitsubject{Quiz on version control basics}
\setreportkitkeywords{ReportKit, quiz, true false, multiple choice}
\setreportkitlanguage{en-US}
\title{Unit 3 Quiz: Version Control Basics}
\author{Fictional Course Staff}
\date{November 2026}

\begin{document}
\maketitle
\quizheader

\quizsection{Part A. True or False}{Circle \emph{True} or \emph{False} for each statement.}

\quiztf{A commit records a snapshot of the tracked files.}{T}
  [A commit stores a snapshot plus metadata such as author and message.]
\quiztf{Deleting a branch always deletes the commits on it.}{F}
  [Commits stay reachable from other refs, and are recoverable for a time otherwise.]
\quiztf{A merge conflict means the repository is corrupted.}{F}
  [It means two changes touched the same lines and need a human decision.]
\quiztf{Every clone of a distributed repository holds the full history.}{T}

\quizsection{Part B. Multiple Choice}{Circle the letter of the single best answer.}

\begin{quizmc}{C}{Which command shows uncommitted changes in the working tree?}
  \quizopt{\texttt{git log}}
  \quizopt{\texttt{git push}}
  \quizopt{\texttt{git status}}
  \quizopt{\texttt{git tag}}
\end{quizmc}
\quizexplain{\texttt{git status} lists staged, unstaged, and untracked files.}

\begin{quizmc}{B}{What is the purpose of a branch?}
  \quizopt{To compress the repository.}
  \quizopt{To isolate a line of work from other work.}
  \quizopt{To encrypt commits before sharing.}
  \quizopt{To remove files from history.}
  \quizopt{To rename the default remote.}
\end{quizmc}

\begin{quizmc}{A}{Which statement about \texttt{git pull} is most accurate?}
  \quizopt{It fetches remote changes and integrates them into the current branch.}
  \quizopt{It uploads local commits to the remote.}
  \quizopt{It creates a new branch from a tag.}
  \quizopt{It discards all local changes.}
\end{quizmc}
\quizexplain{Pull is a fetch followed by a merge (or rebase).}

\vspace{1.5\baselineskip}\noindent\emph{Total questions: \quizscore.}

\end{document}
